\scanpage{85}
% archive sheet number: 84
\unsure{Et} la matrice de $\sigma_1$ est.
\[
\sigma_1=\begin{pmatrix}0&1&-(1+\alpha)&0\\1&0&1+\alpha&0\\0&0&1&0\\0&0&0&1\end{pmatrix}
=\begin{pmatrix}0&1&\alpha'&0\\1&0&-\alpha'&0\\0&0&1&0\\0&0&0&1\end{pmatrix}\tag*{(3)}
\]
Posons
\[
\left\{\begin{aligned}
u_{\deleted{\unsure{0|1}}}&=\sigma_0\sigma_1
& s_0\mapsto s_1\mapsto s_2\mapsto\unsure{s_1|s_0},\ s_3\mapsto s''_3\deleted{\unsure{4}}\\
v_{\deleted{\unsure{1|2}}}&=\sigma_1\sigma_2
& s_0\mapsto s_0,\ s_3\mapsto s_1\mapsto s_2\mapsto s'_3
\end{aligned}\right.\tag*{(4)}
\]
\margin{à calculer $\unsure{\nwarrow}$}
\deleted{Notons que $\unsure{s_0},s_1,s_2,s_3$}

On aura
\[
\begin{aligned}
s''_3&=us_0+ue_3=s_{\unsure{1|2}}+\sigma_0(\sigma_1(e_3))=s_0+e_1+\sigma_0\bigl[(1+\alpha)(e_2-e_1)\\
&\hspace{10cm}+e_3\bigr]\\
&=s_0+e_1+(1+\alpha)\deleted{\unsure{(}}e_2\deleted{\unsure{-e_1)}}+e_3-e_1
\end{aligned}
\]
donc
\[
s''_3=s_0+(1+\alpha)\deleted{\unsure{(}}e_2\deleted{\unsure{-e_1)}}+e_3=s_3+(1+\alpha)e_2\tag*{(5)}
\]
Ou encore
\[
s''_3-s_3=(1+\alpha)(s_2-s_0)\tag*{(6)}
\]
\unsure{Notons|Calculons} $u(s'_3)$ utilisant (2)
\[
\deleted{us'_3=us_0+ue_3=s_1+\sigma_0\sigma_1\bigl((1+\alpha)(e_2-e_1)+e_3\bigr)=s_1+(1+\alpha)u(s_2-s_1}
\]
\[
\begin{aligned}
u(s'_3)&=u(s_3)+(1+\alpha)\underbrace{u(s_2-s_1)}_{s_0-s_2=-e_2}=s''_3-(1+\alpha)e_2\\
&\underset{\unsure{(5)|(6)}}{=}s_3
\end{aligned}
\]
\deleted{\unsure{au}} Enfin, on trouve de $\unsure{=|\text{même}}$ par (6),
\[
\deleted{\unsure{(7)}}\qquad\deleted{\unsure{u(s_3)\quad s_3\longrightarrow s'}}
\]
\[
\begin{aligned}
\deleted{\unsure{(7)}}\quad u(s''_3)&=u(s_3)+(1+\alpha)\overbrace{u(s_2-s_0)}^{s_0-s_1=-e_1}=s''_3-(1+\alpha)e_1\\
&=s_0+(1+\alpha)(e_2-e_1)+e_3
\end{aligned}
\]
